% Options for packages loaded elsewhere
\PassOptionsToPackage{unicode}{hyperref}
\PassOptionsToPackage{hyphens}{url}
\PassOptionsToPackage{dvipsnames,svgnames,x11names}{xcolor}
%
\documentclass[
  10pt,
]{article}
\usepackage{amsmath,amssymb}
\usepackage{iftex}
\ifPDFTeX
  \usepackage[T1]{fontenc}
  \usepackage[utf8]{inputenc}
  \usepackage{textcomp} % provide euro and other symbols
\else % if luatex or xetex
  \usepackage{unicode-math} % this also loads fontspec
  \defaultfontfeatures{Scale=MatchLowercase}
  \defaultfontfeatures[\rmfamily]{Ligatures=TeX,Scale=1}
\fi
\usepackage{lmodern}
\ifPDFTeX\else
  % xetex/luatex font selection
  \setmainfont[]{DejaVu Serif}
  \setsansfont[]{DejaVu Sans}
  \setmonofont[]{DejaVu Sans Mono}
  \setmathfont[]{Latin Modern Math}
\fi
% Use upquote if available, for straight quotes in verbatim environments
\IfFileExists{upquote.sty}{\usepackage{upquote}}{}
\IfFileExists{microtype.sty}{% use microtype if available
  \usepackage[]{microtype}
  \UseMicrotypeSet[protrusion]{basicmath} % disable protrusion for tt fonts
}{}
\makeatletter
\@ifundefined{KOMAClassName}{% if non-KOMA class
  \IfFileExists{parskip.sty}{%
    \usepackage{parskip}
  }{% else
    \setlength{\parindent}{0pt}
    \setlength{\parskip}{6pt plus 2pt minus 1pt}}
}{% if KOMA class
  \KOMAoptions{parskip=half}}
\makeatother
\usepackage{xcolor}
\usepackage[margin=2.0cm]{geometry}
\usepackage{color}
\usepackage{fancyvrb}
\newcommand{\VerbBar}{|}
\newcommand{\VERB}{\Verb[commandchars=\\\{\}]}
\DefineVerbatimEnvironment{Highlighting}{Verbatim}{commandchars=\\\{\}}
% Add ',fontsize=\small' for more characters per line
\newenvironment{Shaded}{}{}
\newcommand{\AlertTok}[1]{\textcolor[rgb]{1.00,0.00,0.00}{\textbf{#1}}}
\newcommand{\AnnotationTok}[1]{\textcolor[rgb]{0.38,0.63,0.69}{\textbf{\textit{#1}}}}
\newcommand{\AttributeTok}[1]{\textcolor[rgb]{0.49,0.56,0.16}{#1}}
\newcommand{\BaseNTok}[1]{\textcolor[rgb]{0.25,0.63,0.44}{#1}}
\newcommand{\BuiltInTok}[1]{\textcolor[rgb]{0.00,0.50,0.00}{#1}}
\newcommand{\CharTok}[1]{\textcolor[rgb]{0.25,0.44,0.63}{#1}}
\newcommand{\CommentTok}[1]{\textcolor[rgb]{0.38,0.63,0.69}{\textit{#1}}}
\newcommand{\CommentVarTok}[1]{\textcolor[rgb]{0.38,0.63,0.69}{\textbf{\textit{#1}}}}
\newcommand{\ConstantTok}[1]{\textcolor[rgb]{0.53,0.00,0.00}{#1}}
\newcommand{\ControlFlowTok}[1]{\textcolor[rgb]{0.00,0.44,0.13}{\textbf{#1}}}
\newcommand{\DataTypeTok}[1]{\textcolor[rgb]{0.56,0.13,0.00}{#1}}
\newcommand{\DecValTok}[1]{\textcolor[rgb]{0.25,0.63,0.44}{#1}}
\newcommand{\DocumentationTok}[1]{\textcolor[rgb]{0.73,0.13,0.13}{\textit{#1}}}
\newcommand{\ErrorTok}[1]{\textcolor[rgb]{1.00,0.00,0.00}{\textbf{#1}}}
\newcommand{\ExtensionTok}[1]{#1}
\newcommand{\FloatTok}[1]{\textcolor[rgb]{0.25,0.63,0.44}{#1}}
\newcommand{\FunctionTok}[1]{\textcolor[rgb]{0.02,0.16,0.49}{#1}}
\newcommand{\ImportTok}[1]{\textcolor[rgb]{0.00,0.50,0.00}{\textbf{#1}}}
\newcommand{\InformationTok}[1]{\textcolor[rgb]{0.38,0.63,0.69}{\textbf{\textit{#1}}}}
\newcommand{\KeywordTok}[1]{\textcolor[rgb]{0.00,0.44,0.13}{\textbf{#1}}}
\newcommand{\NormalTok}[1]{#1}
\newcommand{\OperatorTok}[1]{\textcolor[rgb]{0.40,0.40,0.40}{#1}}
\newcommand{\OtherTok}[1]{\textcolor[rgb]{0.00,0.44,0.13}{#1}}
\newcommand{\PreprocessorTok}[1]{\textcolor[rgb]{0.74,0.48,0.00}{#1}}
\newcommand{\RegionMarkerTok}[1]{#1}
\newcommand{\SpecialCharTok}[1]{\textcolor[rgb]{0.25,0.44,0.63}{#1}}
\newcommand{\SpecialStringTok}[1]{\textcolor[rgb]{0.73,0.40,0.53}{#1}}
\newcommand{\StringTok}[1]{\textcolor[rgb]{0.25,0.44,0.63}{#1}}
\newcommand{\VariableTok}[1]{\textcolor[rgb]{0.10,0.09,0.49}{#1}}
\newcommand{\VerbatimStringTok}[1]{\textcolor[rgb]{0.25,0.44,0.63}{#1}}
\newcommand{\WarningTok}[1]{\textcolor[rgb]{0.38,0.63,0.69}{\textbf{\textit{#1}}}}
\usepackage{longtable,booktabs,array}
\usepackage{calc} % for calculating minipage widths
% Correct order of tables after \paragraph or \subparagraph
\usepackage{etoolbox}
\makeatletter
\patchcmd\longtable{\par}{\if@noskipsec\mbox{}\fi\par}{}{}
\makeatother
% Allow footnotes in longtable head/foot
\IfFileExists{footnotehyper.sty}{\usepackage{footnotehyper}}{\usepackage{footnote}}
\makesavenoteenv{longtable}
\setlength{\emergencystretch}{3em} % prevent overfull lines
\providecommand{\tightlist}{%
  \setlength{\itemsep}{0pt}\setlength{\parskip}{0pt}}
\setcounter{secnumdepth}{-\maxdimen} % remove section numbering
% doc_header.tex -- LaTeX preamble additions used by build_docs.sh for every
% Markdown -> LaTeX -> PDF conversion (supplements, README, submission files).
\usepackage{etoolbox}
% longtable may not be loaded by pandoc when a document has no table
\usepackage{longtable}
% bold vectors: \bm -> \symbf under unicode-math (XeLaTeX), bm package otherwise
\makeatletter
\@ifpackageloaded{unicode-math}{\providecommand{\bm}[1]{\symbf{#1}}}{\usepackage{bm}}
\makeatother
% smaller type in long tables so that wide audit tables fit the text width;
% inside a table a line break is allowed after an underscore of plain text
% (pandoc writes it as \_), so that long identifiers wrap inside their column
% rows are set a little apart, so that symbols with indices do not touch
\AtBeginEnvironment{longtable}{\footnotesize\renewcommand{\arraystretch}{1.45}%
  \let\sptunderscore\_\def\_{\sptunderscore\allowbreak}}
\setlength{\LTcapwidth}{\textwidth}
% allow line breaks after slashes and hyphens inside table cells
\usepackage{xurl}
% code blocks: break over-long lines instead of running into the margin.
% Plain ``` blocks are typeset by pandoc as {verbatim}, highlighted ones
% (```bash) as {Highlighting}; both are re-declared as breakable fancyvrb
% environments.  The continuation line is marked with a small hooked arrow.
\usepackage{fvextra}
\RecustomVerbatimEnvironment{verbatim}{Verbatim}{breaklines,breakanywhere,fontsize=\small}
\makeatletter
\@ifundefined{Highlighting}{}{%
  \RecustomVerbatimEnvironment{Highlighting}{Verbatim}{breaklines,breakanywhere,fontsize=\small,commandchars=\\\{\}}}
\makeatother
% inline code (pandoc writes it as \texttt{...}): never hyphenate -- an inserted
% hyphen could be read as part of a file name -- but allow a line break after
% an underscore, slash, comma, colon or equals sign, so that long paths and
% column lists wrap inside the text width and inside table columns.
% Inside a formula (e.g. the name of a notebook symbol in an equation) the same
% command must first leave math mode, which \text of amsmath does.
\ExplSyntaxOn
\tl_new:N \l__spt_code_tl
\cs_new_protected:Npn \__spt_texttt:n #1
  {
    \tl_set:Nn \l__spt_code_tl {#1}
    \regex_replace_all:nnN { ([\/\,\:\=]) } { \1 \c{allowbreak} } \l__spt_code_tl
    \regex_replace_all:nnN { \c{_} } { \c{_} \c{allowbreak} } \l__spt_code_tl
    \group_begin: \ttfamily \hyphenchar \font = -1 \scan_stop: \l__spt_code_tl \group_end:
  }
\RenewDocumentCommand \texttt { m }
  {
    \mode_if_math:TF { \text { \__spt_texttt:n {#1} } } { \__spt_texttt:n {#1} }
  }
\ExplSyntaxOff
% A figure is written in the Markdown source as an image paragraph followed by
% a caption paragraph ("**Figure S1.** ..."), so that the caption is also
% visible on GitHub.  An image that starts a paragraph is centered.
\usepackage{graphicx}% pandoc loads it only for documents that contain an image
\NewCommandCopy\sptincludegraphics\includegraphics
\RenewDocumentCommand\includegraphics{O{} m}{%
  \ifvmode\begin{center}\sptincludegraphics[#1]{#2}\end{center}%
  \else\sptincludegraphics[#1]{#2}\fi}
% keep figures where they are referenced
\usepackage{float}
\floatplacement{figure}{H}
% running header with the document title
\usepackage{fancyhdr}
\pagestyle{fancy}
\fancyhf{}
\fancyhead[L]{\small PROJECT spacetimes, Revision\_Deflation --- documentation and provenance}
\fancyhead[R]{\small\thepage}
\renewcommand{\headrulewidth}{0.4pt}
% 16x16 matrices: tighter columns so that a full matrix fits the text width
\setlength{\arraycolsep}{3.2pt}
% The documents of this project contain long tagged symbols, file paths and
% identifiers that cannot be hyphenated.  A paragraph is therefore allowed to
% have looser lines before TeX gives up and lets a line stick out of the margin.
\AtBeginDocument{\tolerance=2000 \emergencystretch=5em \hbadness=2000}
\ifLuaTeX
  \usepackage{selnolig}  % disable illegal ligatures
\fi
\IfFileExists{bookmark.sty}{\usepackage{bookmark}}{\usepackage{hyperref}}
\IfFileExists{xurl.sty}{\usepackage{xurl}}{} % add URL line breaks if available
\urlstyle{same}
\hypersetup{
  pdftitle={RP0 --- The hypotheses of Revision\_Deflation, the assumptions added to them, what was repeated, and what is open},
  pdfauthor={Prepared for Patrick Lee Nash, Ph.D.},
  colorlinks=true,
  linkcolor={blue},
  filecolor={Maroon},
  citecolor={blue},
  urlcolor={blue},
  pdfcreator={LaTeX via pandoc}}

\title{RP0 --- The hypotheses of Revision\_Deflation, the assumptions
added to them, what was repeated, and what is open}
\usepackage{etoolbox}
\makeatletter
\providecommand{\subtitle}[1]{% add subtitle to \maketitle
  \apptocmd{\@title}{\par {\large #1 \par}}{}{}
}
\makeatother
\subtitle{Provenance of PROJECT spacetimes, Revision\_Deflation: the map
of the revision, and the account required by Hard Rule 7 of the one
calculation that needed an assumption beyond the hypotheses}
\author{Prepared for Patrick Lee Nash, Ph.D.}
\date{November 13, 2026}

\begin{document}
\maketitle

\hypertarget{the-hypotheses}{%
\section{1. The hypotheses}\label{the-hypotheses}}

On November 12, 2026 the author paused all stages of the main project
(they are recorded in the provenance P11 of the main project and are to
be completed later) and instructed that every calculation of the main
task be repeated, in the new folder \texttt{Revision\_Deflation}, under
new hypotheses. In the assistant's words, with the author's formulas:

\begin{itemize}
\tightlist
\item
  On the 4+4 spacetime and on the 7+1 spacetime alike, the coordinates
  \(x5\), \(x6\) and \(x7\) \textbf{deflate exponentially}.
\item
  Minkowski space is replaced by a pseudo-Riemannian space with a
  \textbf{diagonal metric} whose three components along these
  coordinates are, for the 4+4 spacetime,
  \[g_{55}=-\mathrm{const}^2\,e^{-2\,\mathrm{g55}\,x4},\qquad
  g_{66}=-\mathrm{const}^2\,e^{-2\,\mathrm{g66}\,x4},\qquad
  g_{77}=-\mathrm{const}^2\,e^{-2\,\mathrm{g77}\,x4},\] with
  \(\mathrm{g55},\mathrm{g66},\mathrm{g77}>0\) and constants that can be
  different. (The author first wrote the three components with a
  positive sign and then corrected the sign for the 4+4 spacetime.)
\item
  Everything that the main task demands is to be mirrored, including two
  PRD manuscripts with the titles ``16 Component Dirac Fermion
  Condensate Singlet Superfluid on 4+4 Deflating Spacetime'' and
  ``\ldots{} on 7+1 Deflating Spacetime''.
\end{itemize}

\textbf{How the hypotheses were read.}

\begin{enumerate}
\def\labelenumi{\arabic{enumi}.}
\tightlist
\item
  \emph{The names.} \(\mathrm{g55}\), \(\mathrm{g66}\), \(\mathrm{g77}\)
  are the author's names for the three rates, and they are used exactly
  so, upright and without subscript. The components of the metric are
  written \(g_{55}\), \(g_{66}\), \(g_{77}\). (The assistant at first
  replaced the author's names by a Greek letter, which the author
  rejected; error R-E4 of the provenance RP9.) The author's word const
  stands for three constants, which are named \(\mathrm{const5}\),
  \(\mathrm{const6}\), \(\mathrm{const7}\).
\item
  \emph{The sign on 7+1.} The positive sign of the author's first
  formula is the sign of these components of the metric in the
  convention in which space-like components are positive. The main
  project writes the 7+1 metric with the time-like component positive,
  \(\eta_{1671}=\mathrm{diag}(-1,-1,-1,-1,+1,-1,-1,-1)\). In that
  convention the three components are negative on 7+1 as well. The
  geometry --- a proper length along \(x5\) that shrinks as
  \(e^{-\mathrm{g55}\,x4}\) --- is the same in either convention.
\item
  \emph{The other coordinates} keep their flat metric: \(x0\) (the
  hidden space coordinate), \(x1,x2,x3\) (space) and \(x4\) (the
  ordinary time).
\item
  \emph{``Fresh calculations''.} Nothing was taken over from the main
  project without being computed again on the deflating spacetimes;
  where a result of the main project holds unchanged, the program of the
  revision verifies that it does (for example, all 106 checks of the
  Dirac matrices, the 1170 exchange amplitudes, the internal symmetry).
\end{enumerate}

\hypertarget{what-the-revision-contains}{%
\section{2. What the revision
contains}\label{what-the-revision-contains}}

\begin{longtable}[]{@{}
  >{\raggedright\arraybackslash}p{(\columnwidth - 2\tabcolsep) * \real{0.2353}}
  >{\raggedright\arraybackslash}p{(\columnwidth - 2\tabcolsep) * \real{0.7647}}@{}}
\toprule\noalign{}
\begin{minipage}[b]{\linewidth}\raggedright
document
\end{minipage} & \begin{minipage}[b]{\linewidth}\raggedright
content
\end{minipage} \\
\midrule\noalign{}
\endhead
\bottomrule\noalign{}
\endlastfoot
\texttt{dirac\ matrices\ deflating} (DMD) & geometry of the deflating
spacetimes; Dirac matrices, spin connection; the Dirac equations of
dirac1644d and dirac1671d; their solutions; symmetries \\
\texttt{supplement/RS1} & the Vorton-Scalar and the superfluid:
reduction to four dimensions, the deflation clock, pairing, gap, bound
pair, order parameter, vortices, dynamics \\
\texttt{supplement/RS2} & the Higgs hypothesis H2 on the deflating
spacetimes \\
\texttt{supplement/RS3} & cosmology: bounds on the rate, Hubble
constant, dark matter \\
\texttt{supplement/RS4} & the effective action of Cornwall, Jackiw and
Tomboulis along the clock \\
\texttt{supplement/RS5} & the microwave background and the growth of
structure: perturbation equations, closures, spectra, forecasts \\
\texttt{provenance/RP0} & this document \\
\texttt{provenance/RP1} & MASTER provenance of the symbols of the
revision (Hard Rule 5) \\
\texttt{provenance/RP2}, \texttt{RP3} & the PDF files and the notebooks
used \\
\texttt{provenance/RP8} & the source code: what computes what, how to
compile and run it \\
\texttt{provenance/RP9} & the error log of the revision \\
\texttt{manuscript/} & the two PRD manuscripts \\
\texttt{README} & where to start \\
\end{longtable}

Every document exists as \texttt{.md}, \texttt{.tex} and \texttt{.pdf}
(Hard Rule 1).

\hypertarget{the-assumptions-that-were-added-and-the-hypotheses-that-were-tested}{%
\section{3. The assumptions that were added, and the hypotheses that
were
tested}\label{the-assumptions-that-were-added-and-the-hypotheses-that-were-tested}}

The author's hypotheses fix the metric. Four further assumptions were
needed in the cosmological part, and one hypothesis of the main project
was tested again. They are kept apart in every document.

\begin{longtable}[]{@{}
  >{\raggedright\arraybackslash}p{(\columnwidth - 6\tabcolsep) * \real{0.1111}}
  >{\raggedright\arraybackslash}p{(\columnwidth - 6\tabcolsep) * \real{0.5926}}
  >{\raggedright\arraybackslash}p{(\columnwidth - 6\tabcolsep) * \real{0.1481}}
  >{\raggedright\arraybackslash}p{(\columnwidth - 6\tabcolsep) * \real{0.1481}}@{}}
\toprule\noalign{}
\begin{minipage}[b]{\linewidth}\raggedright
label
\end{minipage} & \begin{minipage}[b]{\linewidth}\raggedright
statement
\end{minipage} & \begin{minipage}[b]{\linewidth}\raggedright
whose
\end{minipage} & \begin{minipage}[b]{\linewidth}\raggedright
where used
\end{minipage} \\
\midrule\noalign{}
\endhead
\bottomrule\noalign{}
\endlastfoot
A1 & ordinary space \(x1,x2,x3\) expands with the scale factor of
standard cosmology & assumption of the revision & RS3 \\
A2 & gravity is the Einstein theory in eight dimensions, so that
Newton's constant is \(G_8\) divided by the volume of the extra
coordinates and grows at the rate \(\mathrm{gsum}\) & assumption of the
revision & RS3 \\
A3 & in a universe with small inhomogeneities (i) the scale factors of
\(x5,x6,x7\) and of \(x0\) are not perturbed, (ii) whatever keeps the
three rates constant (the \emph{agent}) has no energy density, and (iii)
one further statement about the perturbations of the agent holds (a
\emph{closure}: ``Momentum'', ``Pressure'' or ``Shear'') & assumption of
the revision; see Sec. 5 & RS5; RS3, Sec. 5 \\
HV & the vacuum value of the Vorton-Scalar sets the electroweak scale &
assumption of the revision, examined only to show its consequence & RS3,
Sec. 3 \\
H2 & the superfluid state is the Higgs field & the author's hypothesis
(main project) & RS2; RS3, Sec. 3 \\
\end{longtable}

\hypertarget{what-mirrors-what}{%
\section{4. What mirrors what}\label{what-mirrors-what}}

The program of the main project has eleven commands. The table lists
what each computes and where the revision repeats it. The program of the
revision is \texttt{deflation} (folder \texttt{code/deflation-rs}); its
commands are \texttt{matrices}, \texttt{modes}, \texttt{history},
\texttt{quench}, \texttt{cosmology}, \texttt{symmetry},
\texttt{channels}, \texttt{matter}, \texttt{paired},
\texttt{background}, \texttt{spectrum}, \texttt{forecast},
\texttt{families}, \texttt{superfluid}.

\begin{longtable}[]{@{}
  >{\raggedright\arraybackslash}p{(\columnwidth - 6\tabcolsep) * \real{0.1429}}
  >{\raggedright\arraybackslash}p{(\columnwidth - 6\tabcolsep) * \real{0.3571}}
  >{\raggedright\arraybackslash}p{(\columnwidth - 6\tabcolsep) * \real{0.3571}}
  >{\raggedright\arraybackslash}p{(\columnwidth - 6\tabcolsep) * \real{0.1429}}@{}}
\toprule\noalign{}
\begin{minipage}[b]{\linewidth}\raggedright
command of the main project
\end{minipage} & \begin{minipage}[b]{\linewidth}\raggedright
what it computes
\end{minipage} & \begin{minipage}[b]{\linewidth}\raggedright
repeated in the revision by
\end{minipage} & \begin{minipage}[b]{\linewidth}\raggedright
status
\end{minipage} \\
\midrule\noalign{}
\endhead
\bottomrule\noalign{}
\endlastfoot
\texttt{matrices} & the Dirac matrices and their 106 checks; pair
classification & \texttt{deflation\ matrices}: the same 106 checks
again, and 47 of 47 checks of the geometry, of the coordinate matrices,
of the spin connection and of the Dirac operators & done (DMD) \\
\texttt{channels} & exchange amplitudes of the Vorton-Scalar in 15
representative pair channels on each spacetime &
\texttt{deflation\ channels} (checks \texttt{def-K1} to \texttt{def-K3})
& done (RS1, Sec. 3) \\
\texttt{gap} & the gap equation & \texttt{deflation\ history} (checks
\texttt{def-H2}, \texttt{def-H3}) & done (RS1, Sec. 3) \\
\texttt{dimer} & the bound pair & \texttt{deflation\ history}
(\texttt{def-H1}, \texttt{def-H4}), \texttt{deflation\ matter}
(\texttt{def-D3}) & done (RS1, Sec. 4) \\
\texttt{symmetry} & the internal symmetry of the sector &
\texttt{deflation\ symmetry} (\texttt{def-S3}; new: \texttt{44d-S1},
\texttt{S2}, \texttt{71d-S1}, \texttt{S2}) & done (RS2; DMD, Part 5) \\
\texttt{cjt} & the effective action: paired state, matter, temperature,
vacuum, test of the method & \texttt{deflation\ paired} (\texttt{def-P1}
to \texttt{def-P7}), \texttt{deflation\ history} (\texttt{def-H5} to
\texttt{def-H8}) & done (RS4) \\
\texttt{halos} & estimates for dark-matter halos; self-interaction &
\texttt{deflation\ matter} (\texttt{def-D1}, \texttt{def-D2}) & done
(RS3, Sec. 6) \\
\texttt{background} & the cosmological background and the thermal
history: sound horizon, angular scale, last scattering, helium &
\texttt{deflation\ cosmology} (\texttt{def-C1} to \texttt{def-C4}),
\texttt{deflation\ background} (\texttt{def-B1}, \texttt{def-B2}),
independently checked & done (RS3, Secs. 1 to 4; RS5, Secs. 5, 6) \\
\texttt{superfluid} & the superfluid as a cosmological fluid:
suppression of structure by first sound, limits from the
Lyman-\(\alpha\) forest & \texttt{deflation\ superfluid}
(\texttt{def-U1}, \texttt{def-U2}); in addition
\texttt{deflation\ matter}: growth of structure (\texttt{def-D4}), table
of drifts (\texttt{def-D5}) & done (RS5, Sec. 10; RS3, Secs. 5, 6) \\
\texttt{spectrum} & spectra of the microwave background with the
Einstein--Boltzmann program of the main project &
\texttt{deflation\ spectrum} (\texttt{def-E1} to \texttt{def-E8}), with
copies of four modules of that program & done under A3 (RS5, Secs. 2 to
7) \\
\texttt{forecast} & forecast of constraints from mock data, for seven
families of dark radiation & \texttt{deflation\ forecast}
(\texttt{def-F1} to \texttt{def-F5}: the rate of deflation as a
parameter) and \texttt{deflation\ families} (\texttt{def-N1} to
\texttt{def-N4}: the seven families on the deflating background, and
together with the rate) & done under A3 (RS5, Secs. 8, 9) \\
\end{longtable}

New in the revision, without counterpart in the main project:

\begin{longtable}[]{@{}
  >{\raggedright\arraybackslash}p{(\columnwidth - 4\tabcolsep) * \real{0.2222}}
  >{\raggedright\arraybackslash}p{(\columnwidth - 4\tabcolsep) * \real{0.5556}}
  >{\raggedright\arraybackslash}p{(\columnwidth - 4\tabcolsep) * \real{0.2222}}@{}}
\toprule\noalign{}
\begin{minipage}[b]{\linewidth}\raggedright
command
\end{minipage} & \begin{minipage}[b]{\linewidth}\raggedright
what it computes
\end{minipage} & \begin{minipage}[b]{\linewidth}\raggedright
where
\end{minipage} \\
\midrule\noalign{}
\endhead
\bottomrule\noalign{}
\endlastfoot
\texttt{deflation\ matrices} & geometry: curvature, Einstein and Lanczos
tensors, Killing vectors, geodesics & DMD, Part 1 \\
\texttt{deflation\ modes} & exact and numerical solutions of the Dirac
equations: creation of particles (7+1), extinction (4+4) & DMD, Part
4 \\
\texttt{deflation\ channels} & modes of the scalar: static source,
creation, bounded solutions & RS1, Sec. 2 \\
\texttt{deflation\ symmetry} & which rotations and boosts survive;
vortex zero modes; reduction of the couplings & DMD, Part 5; RS1, Secs.
1, 6 \\
\texttt{deflation\ quench} & whether the paired state can follow the
deflation & RS1, Sec. 7 \\
\texttt{deflation\ cosmology} & bounds on the rate of deflation; the
exponent of the electroweak scale under H2 & RS3, Sec. 3; RS2, Sec. 5 \\
\end{longtable}

In total the program of the revision performs 238 checks, and the
independent Python script 43 of 43; all pass.

The documents of the main project and their counterparts:

\begin{longtable}[]{@{}
  >{\raggedright\arraybackslash}p{(\columnwidth - 2\tabcolsep) * \real{0.5000}}
  >{\raggedright\arraybackslash}p{(\columnwidth - 2\tabcolsep) * \real{0.5000}}@{}}
\toprule\noalign{}
\begin{minipage}[b]{\linewidth}\raggedright
main project
\end{minipage} & \begin{minipage}[b]{\linewidth}\raggedright
revision
\end{minipage} \\
\midrule\noalign{}
\endhead
\bottomrule\noalign{}
\endlastfoot
\texttt{dirac\ matrices} (DM) & \texttt{dirac\ matrices\ deflating}
(DMD) \\
supplements S1 to S4 & supplements RS1 to RS5 (S3 is mirrored by RS3 and
RS5) \\
the two manuscripts on the flat 4+4 and 7+1 spacetimes & the two
manuscripts on the deflating 4+4 and 7+1 spacetimes (folder
\texttt{manuscript}), built from one source
(\texttt{docs\_src/manuscript}) by
\texttt{code/scripts/build\_manuscripts.py} \\
provenances P0 to P3, P8, P9 & provenances RP0 to RP3, RP8, RP9 \\
provenance P11 (paused stages) & none: nothing of the revision is
paused; Sec. 5 below lists what the author is asked to confirm \\
\end{longtable}

\hypertarget{the-calculation-that-needed-an-assumption-beyond-the-hypotheses-hard-rule-7}{%
\section{5. The calculation that needed an assumption beyond the
hypotheses (Hard Rule
7)}\label{the-calculation-that-needed-an-assumption-beyond-the-hypotheses-hard-rule-7}}

\textbf{What.} The commands \texttt{spectrum}, \texttt{forecast} and
\texttt{superfluid} of the main project compute the spectra of the
microwave background and the power spectrum of matter with an
Einstein--Boltzmann program written for the main project.

\textbf{Why the hypotheses do not determine it.} The unperturbed
universe of RS3 needs one of Einstein's equations, the component along
the time. The spectra are built from small perturbations of ordinary
space, and these obey the other components as well, which contain the
stresses of whatever keeps the three rates of deflation constant. The
hypotheses of the revision prescribe the metric and say nothing about
that agent. An earlier version of this section therefore recorded the
three calculations as \emph{not carried out}, named the simplest
candidate for the missing assumption, and asked the author how to
proceed, as Hard Rule 7 prescribes.

\textbf{What was done then.} The author's reply was the instruction to
continue; it did not choose an assumption. The calculations were
therefore carried out under the candidate that had been named, now
called \textbf{assumption A3} and labelled as an assumption of the
revision wherever it is used, so that the author can accept it or
replace it:

\begin{itemize}
\tightlist
\item
  \emph{(i)} the scale factors of \(x5\), \(x6\), \(x7\) (and of \(x0\))
  are not perturbed: they remain the prescribed functions of the time in
  a universe with small inhomogeneities, the time being the proper time
  of observers who move with the cold matter;
\item
  \emph{(ii)} the agent has no energy density, in the perturbed universe
  as in the unperturbed one (this is assumption A2 taken to first
  order);
\item
  \emph{(iii)} a closure.
\end{itemize}

Part (iii) was not foreseen. The earlier version of this section
expected that, with unperturbed extra scale factors, ``the perturbations
of ordinary spacetime obey the four-dimensional equations with the
Newton's constant of the moment''. The derivation (RS5, Secs. 2, 3;
computer algebra) shows that this is not a consistent statement: of the
four scalar Einstein equations of ordinary space, which in four
dimensions imply one another, only two can hold with matter and
radiation alone when Newton's constant depends on the time. The agent
must carry perturbations of its momentum density, its pressure or its
anisotropic stress, of the relative size \(\mathrm{gsum}/H\), and one
must say which. The revision does not choose. It computes with two
closures (``Momentum'': no momentum density; ``Pressure'': no pressure
perturbation), shows analytically and numerically that a third
(``Shear'': no anisotropic stress) coincides with ``Pressure'' on scales
far inside the horizon, and marks every result as independent of the
closure or as depending on it:

\begin{longtable}[]{@{}
  >{\raggedright\arraybackslash}p{(\columnwidth - 2\tabcolsep) * \real{0.7647}}
  >{\raggedright\arraybackslash}p{(\columnwidth - 2\tabcolsep) * \real{0.2353}}@{}}
\toprule\noalign{}
\begin{minipage}[b]{\linewidth}\raggedright
result
\end{minipage} & \begin{minipage}[b]{\linewidth}\raggedright
depends on the closure?
\end{minipage} \\
\midrule\noalign{}
\endhead
\bottomrule\noalign{}
\endlastfoot
the Hubble constant at a fixed angle of the sound horizon;
recombination; helium & no (these need A1 and A2 only) \\
the spectra for \(\ell\ge30\); the forecast for the rate of deflation;
the fits of the dark-radiation families & no (the closures agree to a
thousandth of the effect of the deflation) \\
the suppression of structure by the first sound of superfluid dark
matter & no (to \(1.5\times10^{-5}\)) \\
the growth of structure, the amplitude \(\sigma_8\) & \textbf{yes}, at
the first order in \(\mathrm{gsum}/H_0\) \\
the spectra for \(\ell<30\) & \textbf{yes} \\
\end{longtable}

\textbf{How it was done without changing the main project.} The main
project is paused, and its program was not changed. The revision
contains copies of four of its modules (background and recombination,
perturbations, spectra, fits). The copies are written by a script,
\texttt{code/scripts/vendor\_eb.py}, which applies a listed set of
replacements to the original files, so that they cannot differ from the
originals in any other place. Without deflation the copies reproduce the
spectra of the main project to the last bit (provenance RP8; RS5, Sec.
5).

\textbf{What the author is asked to decide.}

\begin{enumerate}
\def\labelenumi{\arabic{enumi}.}
\tightlist
\item
  Whether A3, part (i), is what he intends. If the extra scale factors
  respond to inhomogeneities, the equations of RS5 would acquire further
  terms.
\item
  Whether he has a physical model of the agent. It would fix the
  closure, and with it the growth of structure and the largest angular
  scales, which the revision can only bracket.
\item
  Whether gravity is the Einstein theory in eight dimensions (A2) or of
  the Einstein--Gauss--Bonnet kind for which the deflating spacetimes
  with equal rates need no source (DMD, Sec. 1.4a). In the second case
  RS3 and RS5 would have to be derived anew.
\end{enumerate}

No document of the revision presents a result that depends on these
choices without naming them.

Two smaller items are open as well and are stated where they arise: the
exponent \(n_{\rm EW}\) of RS2 is known in closed form only for the
scale-invariant case, and the stability of the Einstein--Gauss--Bonnet
solution of DMD, Sec. 1.4a, has not been examined. The forecasts of RS5
use mock data without gravitational lensing, as those of the main
project do; no measured spectrum has been fitted.

\hypertarget{how-the-revision-reaches-github}{%
\section{6. How the revision reaches
GitHub}\label{how-the-revision-reaches-github}}

Hard Rule 2 demands a push as soon as any object is updated. As
explained in the provenance P0 of the main project, the assistant has no
permission to push to \texttt{github.com/once-ere/spacetimes}; every
commit is transferred to the author's Mac as a git bundle within about a
minute and applied to the local repository there, and a loop that the
author runs in his Terminal pushes it. The same procedure was used for
every commit of the revision. The commits of the revision are listed by

\begin{Shaded}
\begin{Highlighting}[]
\BuiltInTok{cd} \StringTok{"/Users/nsh/Developer/github/spacetimes"}
\FunctionTok{git}\NormalTok{ log }\AttributeTok{{-}{-}oneline} \AttributeTok{{-}{-}}\NormalTok{ Revision\_Deflation}
\end{Highlighting}
\end{Shaded}


\end{document}
